% =============================================================================
% APPENDIX X: LIT-AUTHOR: Unknown Researcher Research Integration
% =============================================================================
% Category: LIT
% Language: English
% Version: 1.0 (January 2026)
% Status: Draft
%
% This appendix integrates research papers by Unknown Researcher into the EBF framework.
%
% =============================================================================

\chapter{LIT-AUTHOR: Unknown Researcher Research Integration}
\label{app:x}

\begin{abstract}
This appendix integrates the research contributions of Unknown Researcher into the
Evidence-Based Framework. It documents key papers, core findings, and their
integration with EBF dimensions including complementarity parameters ($\gamma$),
context factors ($\Psi$), and utility dimensions.
\end{abstract}

\begin{tcolorbox}[colback=blue!5!white, colframe=blue!75!black,
    title=Quick Reference: Unknown Researcher Research]

\textbf{Appendix:} X (LIT)

\textbf{Researcher:} Unknown Researcher

\textbf{Core Themes:}
\begin{itemize}[nosep]
\item Behavioral economics foundations
\item Empirical evidence for EBF parameters
\item Policy implications and interventions
\end{itemize}

\textbf{EBF Integration:}
\begin{itemize}[nosep]
\item Complementarity values ($\gamma$) from experimental evidence
\item Context effects ($\Psi$) documented across studies
\item Utility dimension weightings calibrated to findings
\end{itemize}

\textbf{Cross-References:} BBB (CORE-WHERE), K (LIT-FEHR), G (Glossary)
\end{tcolorbox}


%%%%%%%%%%%%%%%%%%%%%%%%%%%%%%%%%%%%%%%%%%%%%%%%%%%%%%%%%%%%%%%%%%%%%%%%%%%%%%%
\section{The Fundamental Question}
\label{sec:fundamental}
%%%%%%%%%%%%%%%%%%%%%%%%%%%%%%%%%%%%%%%%%%%%%%%%%%%%%%%%%%%%%%%%%%%%%%%%%%%%%%%

\begin{tcolorbox}[colback=red!5!white, colframe=red!75!black,
    title=The Fundamental Question]
What are the key mechanisms, predictions, and implications of this analysis
for the Evidence-Based Framework?
\end{tcolorbox}

% -----------------------------------------------------------------------------
% APPENDIX SCOPE BOX
% -----------------------------------------------------------------------------

\begin{tcolorbox}[
    colback=orange!5!white,
    colframe=orange!75!black,
    title={\textbf{Appendix Scope: Ziel / In-Scope / Out-of-Scope / Constraints}},
    fonttitle=\bfseries
]

\textbf{Ziel (Objective):}
\begin{itemize}[nosep]
    \item Integrate Unknown Researcher's research into EBF with parameter extraction
\end{itemize}

\textbf{In-Scope:}
\begin{itemize}[nosep]
    \item Paper-by-paper integration with 10C mapping
    \item Extraction of $\gamma$, $\Psi$, and effect size parameters
    \item Cross-references to related EBF components
\end{itemize}

\textbf{Out-of-Scope (delegiert an andere Appendices):}
\begin{itemize}[nosep]
    \item Parameter validation $\rightarrow$ Appendix BBB (CORE-WHERE)
    \item Formal axiomatization $\rightarrow$ CORE appendices
    \item Meta-analysis $\rightarrow$ LIT-M appendices
\end{itemize}

\textbf{Constraints (Anwendungsgrenzen):}
\begin{itemize}[nosep]
    \item Parameters are extracted from published studies
    \item Effect sizes may vary across replications
    \item Context-specificity of findings applies
\end{itemize}

\textbf{Lieferobjekte (Deliverables):}
\begin{itemize}[nosep]
    \item Paper integration table with 10C mappings
    \item Extracted parameters for BBB repository
    \item Cross-reference map to related literature
\end{itemize}

\end{tcolorbox}


% =============================================================================
% Appendix VVV7: McKinsey State of AI 2025 Report
% Empirical Evidence for Complementarity-Context Framework
% Source: QuantumBlack, AI by McKinsey (November 2025)
% =============================================================================

% Category: METHOD
% Language: English
\section{McKinsey State of AI 2025: Agents, Innovation, and Transformation}
\label{appendix:VVV1_mckinsey_ai_2025}

\begin{tcolorbox}[colback=blue!5!white, colframe=blue!75!black,
    title=Quick Reference: VVV1 McKinsey Report]
\textbf{Appendix:} ? (METHOD)


\textbf{Source:} McKinsey Global Survey on the State of AI (November 2025)

\textbf{Sample:} n = 1,993 participants, 105 nations

\textbf{Key Finding:} 88\% use AI, but only 38\% are scaling enterprise-wide

\textbf{EBF Relevance:} Empirical validation of institutional complementarity requirements

\end{tcolorbox}


%%%%%%%%%%%%%%%%%%%%%%%%%%%%%%%%%%%%%%%%%%%%%%%%%%%%%%%%%%%%%%%%%%%%%%%%%%%%%%%
\subsection{Overview and Relevance to EBF}
%%%%%%%%%%%%%%%%%%%%%%%%%%%%%%%%%%%%%%%%%%%%%%%%%%%%%%%%%%%%%%%%%%%%%%%%%%%%%%%

This appendix summarizes key findings from McKinsey's Global Survey on the state of AI (November 2025), which surveyed 1,993 participants across 105 nations. The report provides \textbf{empirical evidence} directly relevant to the EBF Complementarity-Context Framework's analysis of:

\begin{itemize}
\item Technological adoption patterns and scaling challenges
\item Institutional complementarities required for AI value capture
\item Context-dependent heterogeneity in AI deployment ($\Psi$-dimensions)
\item The role of leadership and workflow redesign as complementary assets
\end{itemize}


%%%%%%%%%%%%%%%%%%%%%%%%%%%%%%%%%%%%%%%%%%%%%%%%%%%%%%%%%%%%%%%%%%%%%%%%%%%%%%%
\subsection{Executive Summary: Six Key Findings}
%%%%%%%%%%%%%%%%%%%%%%%%%%%%%%%%%%%%%%%%%%%%%%%%%%%%%%%%%%%%%%%%%%%%%%%%%%%%%%%

\begin{enumerate}
\item \textbf{Most organizations remain in experimentation/piloting:} 62\% have not begun scaling AI enterprise-wide
\item \textbf{High curiosity in AI agents:} 62\% are at least experimenting with agentic AI
\item \textbf{Positive leading indicators but limited EBIT impact:} 64\% report innovation benefits, but only 39\% report enterprise EBIT impact
\item \textbf{High performers pursue growth AND innovation:} Not just efficiency---multi-objective optimization matters
\item \textbf{Workflow redesign is critical:} 55\% of high performers fundamentally redesign workflows (vs. 20\% of others)
\item \textbf{Differing workforce expectations:} 32\% expect decreases, 43\% no change, 13\% increases
\end{enumerate}


%%%%%%%%%%%%%%%%%%%%%%%%%%%%%%%%%%%%%%%%%%%%%%%%%%%%%%%%%%%%%%%%%%%%%%%%%%%%%%%
\subsection{AI Adoption Phases}
\label{sec:vvv1-adoption}
%%%%%%%%%%%%%%%%%%%%%%%%%%%%%%%%%%%%%%%%%%%%%%%%%%%%%%%%%%%%%%%%%%%%%%%%%%%%%%%

\begin{table}[htbp]
\centering
\caption{Phase of AI Use Among Organizations (2025)}
\label{tab:vvv1-phases}
\renewcommand{\arraystretch}{1.3}
\begin{tabular}{@{}lcp{6cm}@{}}
\toprule
\textbf{Phase} & \textbf{\%} & \textbf{Definition} \\
\midrule
\textbf{Fully Scaled} & 7\% & AI fully deployed and integrated across organization \\
\textbf{Scaling} & 31\% & Growing deployment/adoption across organization \\
\textbf{Piloting} & 30\% & Implementing AI for first use case in business \\
\textbf{Experimenting} & 32\% & Any use or early testing of AI \\
\midrule
\textit{Total Scaling+} & \textit{38\%} & \textit{Organizations beyond pilot phase} \\
\bottomrule
\end{tabular}
\end{table}

\begin{tcolorbox}[colback=yellow!5!white, colframe=yellow!75!black,
    title=EBF Framework Implication: Institutional Complementarity Gap]

The 62\% still in pilot/experiment phase despite 88\% AI availability demonstrates \textbf{Proposition 4.1 (Institutional Complementarity)}:

\begin{quote}
\textit{``Technological capability alone is insufficient. Value capture requires complementary investments in workflows, talent, leadership, and organizational structure.''}
\end{quote}

This is the \textbf{institutional complementarity gap}---the difference between technology availability and value realization.

\end{tcolorbox}


%%%%%%%%%%%%%%%%%%%%%%%%%%%%%%%%%%%%%%%%%%%%%%%%%%%%%%%%%%%%%%%%%%%%%%%%%%%%%%%
\subsection{Scaling by Company Size}
%%%%%%%%%%%%%%%%%%%%%%%%%%%%%%%%%%%%%%%%%%%%%%%%%%%%%%%%%%%%%%%%%%%%%%%%%%%%%%%

\begin{table}[htbp]
\centering
\caption{AI Scaling Phase by Company Revenue}
\label{tab:vvv1-size}
\renewcommand{\arraystretch}{1.3}
\begin{tabular}{@{}lccccc@{}}
\toprule
\textbf{Revenue} & \textbf{Not Using} & \textbf{Experimenting} & \textbf{Piloting} & \textbf{Scaling} & \textbf{Fully Scaled} \\
\midrule
$<$\$100M & 9\% & 39\% & 22\% & 25\% & 5\% \\
\$100M--\$499M & 8\% & 33\% & 32\% & 23\% & 4\% \\
\$500M--\$999M & 5\% & 31\% & 32\% & 29\% & 3\% \\
\$1B--\$4.9B & 3\% & 22\% & 31\% & 32\% & 9\% \\
\textbf{\$5B+} & \textbf{1\%} & \textbf{17\%} & \textbf{31\%} & \textbf{39\%} & \textbf{10\%} \\
\midrule
\textit{Scaling+ Total} & & & & \multicolumn{2}{c}{30\% $\to$ 49\%} \\
\bottomrule
\end{tabular}
\end{table}

\textbf{Key Insight:} Larger companies (>\$5B revenue) are 1.6× more likely to reach scaling phase (49\%) compared to smaller firms (<\$100M: 30\%). This reflects the \textbf{complementary investment capacity} of larger organizations.


%%%%%%%%%%%%%%%%%%%%%%%%%%%%%%%%%%%%%%%%%%%%%%%%%%%%%%%%%%%%%%%%%%%%%%%%%%%%%%%
\subsection{AI Agents: The Emerging Frontier}
\label{sec:vvv1-agents}
%%%%%%%%%%%%%%%%%%%%%%%%%%%%%%%%%%%%%%%%%%%%%%%%%%%%%%%%%%%%%%%%%%%%%%%%%%%%%%%

AI agents---systems based on foundation models capable of \textit{acting in the real world}, planning and executing multiple steps in a workflow---represent the next wave of AI deployment.

\begin{table}[htbp]
\centering
\caption{AI Agent Scaling by Business Function (\% Scaling or Fully Scaled)}
\label{tab:vvv1-agents}
\renewcommand{\arraystretch}{1.3}
\begin{tabular}{@{}lcc@{}}
\toprule
\textbf{Function} & \textbf{Total} & \textbf{High Performers} \\
\midrule
IT & 10\% & 33\% \\
Knowledge Management & 8\% & 31\% \\
Marketing and Sales & 7\% & 29\% \\
Service Operations & 7\% & 25\% \\
Product Development & 6\% & 29\% \\
Software Engineering & 6\% & 24\% \\
Risk/Legal/Compliance & 6\% & 26\% \\
Human Resources & 6\% & 14\% \\
Strategy \& Finance & 5\% & 21\% \\
Supply Chain & 5\% & 6\% \\
Manufacturing & 4\% & 5\% \\
\bottomrule
\end{tabular}
\end{table}

\begin{tcolorbox}[colback=green!5!white, colframe=green!75!black,
    title=EBF Framework Implication: Agent Adoption = Skills Paradigm]

The agent adoption pattern validates EBF's analysis of the \textbf{GPT $\to$ Skills transition}:

\begin{itemize}
\item \textbf{IT and Knowledge Management lead:} These functions have clear agentic use cases (service desk, deep research)
\item \textbf{High performers are 3-4× ahead:} 33\% vs. 8\% in IT---demonstrating complementarity between ambition and capability
\item \textbf{Manufacturing/Supply Chain lag:} Physical-world integration requires additional complementary infrastructure
\end{itemize}

This heterogeneous adoption pattern reflects \textbf{context-dependent complementarity} ($\Psi$-dimensions):
\begin{equation}
\text{Agent Adoption}_f = g(\text{Use Case Clarity}_f, \text{Complementary Infrastructure}_f, \Psi_{\text{industry}})
\end{equation}

\end{tcolorbox}


%%%%%%%%%%%%%%%%%%%%%%%%%%%%%%%%%%%%%%%%%%%%%%%%%%%%%%%%%%%%%%%%%%%%%%%%%%%%%%%
\subsection{AI Agents by Industry}
%%%%%%%%%%%%%%%%%%%%%%%%%%%%%%%%%%%%%%%%%%%%%%%%%%%%%%%%%%%%%%%%%%%%%%%%%%%%%%%

\begin{table}[htbp]
\centering
\caption{AI Agent Scaling Leaders by Industry}
\label{tab:vvv1-industry}
\renewcommand{\arraystretch}{1.3}
\begin{tabular}{@{}lc@{}}
\toprule
\textbf{Industry} & \textbf{Agent Scaling Rate} \\
\midrule
Technology & Highest \\
Media \& Telecommunications & High \\
Healthcare & High \\
Insurance & High \\
Advanced Manufacturing & Medium \\
Professional Services & Medium \\
Financial Institutions & Medium \\
Consumer Goods \& Retail & Medium \\
Energy \& Materials & Medium \\
Travel \& Logistics & Low \\
Engineering \& Construction & Low \\
Pharma \& Medical Products & Low \\
\bottomrule
\end{tabular}
\end{table}


%%%%%%%%%%%%%%%%%%%%%%%%%%%%%%%%%%%%%%%%%%%%%%%%%%%%%%%%%%%%%%%%%%%%%%%%%%%%%%%
\subsection{Multi-Function AI Deployment Growth}
%%%%%%%%%%%%%%%%%%%%%%%%%%%%%%%%%%%%%%%%%%%%%%%%%%%%%%%%%%%%%%%%%%%%%%%%%%%%%%%

\begin{table}[htbp]
\centering
\caption{AI Use Across Multiple Functions (2021 vs. 2025)}
\label{tab:vvv1-functions}
\renewcommand{\arraystretch}{1.3}
\begin{tabular}{@{}lccc@{}}
\toprule
\textbf{Functions Using AI} & \textbf{2021} & \textbf{2025} & \textbf{Growth} \\
\midrule
1 or more & 56\% & 88\% & +32pp \\
2 or more & 31\% & 70\% & +39pp \\
3 or more & 17\% & 50\% & +33pp \\
4 or more & 9\% & 33\% & +24pp \\
5 or more & 4\% & 20\% & +16pp \\
\bottomrule
\end{tabular}
\end{table}

\textbf{Interpretation:} AI is spreading \textit{horizontally} across organizations---half now use AI in 3+ functions. This creates opportunities for \textbf{cross-functional complementarities} ($\gamma_{ij}$ between AI deployments in different functions).


%%%%%%%%%%%%%%%%%%%%%%%%%%%%%%%%%%%%%%%%%%%%%%%%%%%%%%%%%%%%%%%%%%%%%%%%%%%%%%%
\subsection{Value Creation: Leading Indicators vs. EBIT Impact}
\label{sec:vvv1-value}
%%%%%%%%%%%%%%%%%%%%%%%%%%%%%%%%%%%%%%%%%%%%%%%%%%%%%%%%%%%%%%%%%%%%%%%%%%%%%%%

\begin{table}[htbp]
\centering
\caption{AI Impact on Organizational Measures (\% Reporting Improvement)}
\label{tab:vvv1-impact}
\renewcommand{\arraystretch}{1.3}
\begin{tabular}{@{}lcc@{}}
\toprule
\textbf{Measure} & \textbf{Improved} & \textbf{No Effect} \\
\midrule
\textbf{Innovation} & \textbf{64\%} & 21\% \\
Employee Satisfaction & 45\% & 31\% \\
Customer Satisfaction & 45\% & 32\% \\
Competitive Differentiation & 45\% & 33\% \\
Cost & 38\% & 31\% \\
Profitability & 36\% & 36\% \\
Organic Revenue Growth & 33\% & 39\% \\
Talent Attraction/Retention & 33\% & 42\% \\
Market Share & 25\% & 49\% \\
\midrule
\textit{Enterprise EBIT Impact} & \textit{39\%} & --- \\
\bottomrule
\end{tabular}
\end{table}

\begin{tcolorbox}[colback=red!5!white, colframe=red!75!black,
    title=The Value Realization Gap]

\textbf{64\% report innovation improvement, but only 39\% report EBIT impact.}

This 25 percentage point gap represents the \textbf{value realization challenge}:
\begin{itemize}
\item Organizations are seeing \textit{qualitative} benefits (innovation, satisfaction)
\item But \textit{quantifiable} enterprise-level impact remains elusive
\item This suggests incomplete \textbf{workflow integration} and \textbf{measurement systems}
\end{itemize}

\end{tcolorbox}


%%%%%%%%%%%%%%%%%%%%%%%%%%%%%%%%%%%%%%%%%%%%%%%%%%%%%%%%%%%%%%%%%%%%%%%%%%%%%%%
\subsection{Cost Decreases by Function}
%%%%%%%%%%%%%%%%%%%%%%%%%%%%%%%%%%%%%%%%%%%%%%%%%%%%%%%%%%%%%%%%%%%%%%%%%%%%%%%

\begin{table}[htbp]
\centering
\caption{Cost Decrease from AI Use by Function (Past 12 Months)}
\label{tab:vvv1-cost}
\renewcommand{\arraystretch}{1.3}
\begin{tabular}{@{}lccc@{}}
\toprule
\textbf{Function} & \textbf{$\geq$20\%} & \textbf{11-19\%} & \textbf{Total Decrease} \\
\midrule
Software Engineering & 7\% & 14\% & \textbf{56\%} \\
Manufacturing & 4\% & 7\% & \textbf{56\%} \\
IT & 8\% & 10\% & 54\% \\
Strategy \& Finance & 8\% & 17\% & 53\% \\
Service Operations & 6\% & 8\% & 51\% \\
Human Resources & 7\% & 11\% & 51\% \\
Supply Chain & 7\% & 3\% & 49\% \\
Marketing \& Sales & 7\% & 8\% & 49\% \\
Product Development & 10\% & 10\% & 47\% \\
Risk/Legal/Compliance & 10\% & 11\% & 45\% \\
Knowledge Management & 9\% & 8\% & 44\% \\
\bottomrule
\end{tabular}
\end{table}


%%%%%%%%%%%%%%%%%%%%%%%%%%%%%%%%%%%%%%%%%%%%%%%%%%%%%%%%%%%%%%%%%%%%%%%%%%%%%%%
\subsection{Revenue Increases by Function}
%%%%%%%%%%%%%%%%%%%%%%%%%%%%%%%%%%%%%%%%%%%%%%%%%%%%%%%%%%%%%%%%%%%%%%%%%%%%%%%

\begin{table}[htbp]
\centering
\caption{Revenue Increase from AI Use by Function (Past 12 Months)}
\label{tab:vvv1-revenue}
\renewcommand{\arraystretch}{1.3}
\begin{tabular}{@{}lccc@{}}
\toprule
\textbf{Function} & \textbf{$>$10\%} & \textbf{6-10\%} & \textbf{Total Increase} \\
\midrule
\textbf{Marketing \& Sales} & 10\% & 14\% & \textbf{67\%} \\
\textbf{Strategy \& Finance} & 12\% & 22\% & \textbf{65\%} \\
Product Development & 10\% & 15\% & 62\% \\
Supply Chain & 5\% & 11\% & 59\% \\
Service Operations & 5\% & 16\% & 57\% \\
Software Engineering & 9\% & 13\% & 57\% \\
IT & 5\% & 18\% & 52\% \\
\bottomrule
\end{tabular}
\end{table}


%%%%%%%%%%%%%%%%%%%%%%%%%%%%%%%%%%%%%%%%%%%%%%%%%%%%%%%%%%%%%%%%%%%%%%%%%%%%%%%
\subsection{AI High Performers: Defining Characteristics}
\label{sec:vvv1-high-performers}
%%%%%%%%%%%%%%%%%%%%%%%%%%%%%%%%%%%%%%%%%%%%%%%%%%%%%%%%%%%%%%%%%%%%%%%%%%%%%%%

\textbf{Definition:} AI high performers = respondents reporting $>$5\% EBIT impact from AI AND ``significant value'' (\textbf{6\% of sample}, n=109).

\begin{table}[htbp]
\centering
\caption{High Performers vs. Others: Key Differentiators}
\label{tab:vvv1-hp}
\renewcommand{\arraystretch}{1.4}
\begin{tabular}{@{}lccr@{}}
\toprule
\textbf{Characteristic} & \textbf{High Performers} & \textbf{Others} & \textbf{Ratio} \\
\midrule
\multicolumn{4}{l}{\textit{Ambition \& Strategy}} \\
Aim for transformative change & 50\% & 14\% & \textbf{3.6×} \\
Growth objectives set & 82\% & 50\% & 1.6× \\
Innovation objectives set & 79\% & 50\% & 1.6× \\
Efficiency objectives set & 84\% & 80\% & 1.0× \\
\midrule
\multicolumn{4}{l}{\textit{Execution}} \\
Fundamentally redesign workflows & 55\% & 20\% & \textbf{2.8×} \\
\midrule
\multicolumn{4}{l}{\textit{Leadership}} \\
Leaders demonstrate strong AI commitment & 48\% & 16\% & \textbf{3.0×} \\
\midrule
\multicolumn{4}{l}{\textit{Investment}} \\
$>$20\% of digital budget to AI & 35\% & 7\% & \textbf{4.9×} \\
\bottomrule
\end{tabular}
\end{table}

\begin{tcolorbox}[colback=green!5!white, colframe=green!75!black,
    title=EBF Framework Implication: Supermodularity in Objectives]

High performers don't just pursue efficiency---they pursue \textbf{efficiency + growth + innovation simultaneously}.

This validates \textbf{Proposition 5.3 (Supermodularity)}:
\begin{equation}
\frac{\partial^2 V}{\partial \text{Efficiency} \, \partial \text{Innovation}} > 0
\end{equation}

The value of pursuing efficiency \textit{increases} when also pursuing innovation. The 3.6× differential in transformative ambition creates \textbf{multiplicative, not additive} returns.

\end{tcolorbox}


%%%%%%%%%%%%%%%%%%%%%%%%%%%%%%%%%%%%%%%%%%%%%%%%%%%%%%%%%%%%%%%%%%%%%%%%%%%%%%%
\subsection{Critical Success Factor: Workflow Redesign}
%%%%%%%%%%%%%%%%%%%%%%%%%%%%%%%%%%%%%%%%%%%%%%%%%%%%%%%%%%%%%%%%%%%%%%%%%%%%%%%

The \textbf{2.8× differential} in workflow redesign (55\% vs. 20\%) is identified as ``one of the strongest contributions to achieving meaningful business impact.''

\begin{quote}
\textit{``Leading organizations are not just seeing improved automation results; they are redesigning workflows and customer experiences to capture new forms of value.''} --- Tara Balakrishnan, McKinsey
\end{quote}

\textbf{EBF Interpretation:} AI deployment without workflow redesign is like changing one variable in a complementary system---the system resists. Full value requires \textbf{simultaneous optimization} across technology + process + people.


%%%%%%%%%%%%%%%%%%%%%%%%%%%%%%%%%%%%%%%%%%%%%%%%%%%%%%%%%%%%%%%%%%%%%%%%%%%%%%%
\subsection{Best Practices: The Rewired Framework}
\label{sec:vvv1-practices}
%%%%%%%%%%%%%%%%%%%%%%%%%%%%%%%%%%%%%%%%%%%%%%%%%%%%%%%%%%%%%%%%%%%%%%%%%%%%%%%

McKinsey's ``Rewired'' framework identifies six dimensions essential to AI value capture:

\begin{table}[htbp]
\centering
\caption{Best Practices by Rewired Dimension (\% Adoption)}
\label{tab:vvv1-practices}
\renewcommand{\arraystretch}{1.3}
\begin{tabular}{@{}llcc@{}}
\toprule
\textbf{Dimension} & \textbf{Practice} & \textbf{High Perf.} & \textbf{Others} \\
\midrule
Strategy & Human in the loop processes defined & 65\% & 23\% \\
Strategy & AI vision and strategy clearly defined & 44\% & 21\% \\
Strategy & AI road map with specific initiatives & 60\% & 31\% \\
Strategy & Leaders understand AI value creation & 60\% & 41\% \\
\midrule
Talent & Strategic workforce planning for AI & 54\% & 19\% \\
Talent & AI talent recruitment strategy & 47\% & 18\% \\
Talent & AI upskilling learning journeys & 34\% & 24\% \\
\midrule
Operating Model & Agile product delivery organization & 54\% & 20\% \\
Operating Model & Centralized AI governance team & 46\% & 38\% \\
Operating Model & Rapid development cycles & 54\% & 24\% \\
\midrule
Technology & Infrastructure allows latest AI tech & 60\% & 23\% \\
\midrule
Data & Iterative solution development process & 54\% & 22\% \\
Data & Reusable data products created & 25\% & 21\% \\
\midrule
Adoption & AI embedded in business processes & 58\% & 20\% \\
Adoption & Senior leaders role model AI use & 57\% & 33\% \\
\bottomrule
\end{tabular}
\end{table}

\begin{tcolorbox}[colback=blue!5!white, colframe=blue!75!black,
    title=The Six Complementary Dimensions]

The Rewired framework aligns with EBF's complementarity analysis:

\begin{enumerate}
\item \textbf{Strategy} $\leftrightarrow$ CORE-WHO: Who benefits from AI?
\item \textbf{Talent} $\leftrightarrow$ CORE-AWARE: What capabilities exist?
\item \textbf{Operating Model} $\leftrightarrow$ CORE-HOW: How do components interact?
\item \textbf{Technology} $\leftrightarrow$ CORE-WHERE: What infrastructure exists?
\item \textbf{Data} $\leftrightarrow$ CORE-WHAT: What inputs feed the system?
\item \textbf{Adoption} $\leftrightarrow$ CORE-READY: When do people act?
\end{enumerate}

All six dimensions must be addressed \textit{simultaneously}---they are \textbf{complements, not substitutes}.

\end{tcolorbox}


%%%%%%%%%%%%%%%%%%%%%%%%%%%%%%%%%%%%%%%%%%%%%%%%%%%%%%%%%%%%%%%%%%%%%%%%%%%%%%%
\subsection{Workforce Impact Expectations}
%%%%%%%%%%%%%%%%%%%%%%%%%%%%%%%%%%%%%%%%%%%%%%%%%%%%%%%%%%%%%%%%%%%%%%%%%%%%%%%

\begin{table}[htbp]
\centering
\caption{Expected Change in Enterprise Workforce Size Due to AI (Next Year)}
\label{tab:vvv1-workforce}
\renewcommand{\arraystretch}{1.3}
\begin{tabular}{@{}lc@{}}
\toprule
\textbf{Expected Change} & \textbf{\% of Respondents} \\
\midrule
Decrease by $>$20\% & 4\% \\
Decrease by 11-20\% & 8\% \\
Decrease by 3-10\% & 20\% \\
\textit{Total expecting decrease $\geq$3\%} & \textit{32\%} \\
\midrule
Little or no change & 43\% \\
\midrule
Increase by 3-10\% & 9\% \\
Increase by 11-20\% & 3\% \\
Increase by $>$20\% & 2\% \\
\textit{Total expecting increase $\geq$3\%} & \textit{13\%} \\
\midrule
Don't know & 12\% \\
\bottomrule
\end{tabular}
\end{table}

\textbf{Notable Patterns:}
\begin{itemize}
\item \textbf{Larger organizations} more likely to expect workforce reductions
\item \textbf{High performers} more likely to expect meaningful change \textit{in either direction}
\item This reflects the \textbf{task-level complementarity dynamics} of Section 6.3
\end{itemize}


%%%%%%%%%%%%%%%%%%%%%%%%%%%%%%%%%%%%%%%%%%%%%%%%%%%%%%%%%%%%%%%%%%%%%%%%%%%%%%%
\subsection{AI-Related Hiring}
%%%%%%%%%%%%%%%%%%%%%%%%%%%%%%%%%%%%%%%%%%%%%%%%%%%%%%%%%%%%%%%%%%%%%%%%%%%%%%%

\begin{table}[htbp]
\centering
\caption{AI-Related Hiring in Past Year (Organizations $\geq$\$1B Revenue)}
\label{tab:vvv1-hiring}
\renewcommand{\arraystretch}{1.3}
\begin{tabular}{@{}lc@{}}
\toprule
\textbf{Role} & \textbf{\% Hired} \\
\midrule
AI Data Scientists & 30\% \\
Data Engineers & 29\% \\
Machine Learning Engineers & 29\% \\
Software Engineers & 29\% \\
AI Product Owners/Managers & 26\% \\
Data Architects & 26\% \\
Data Visualization Specialists & 20\% \\
Design Specialists & 18\% \\
AI Compliance Specialists & 16\% \\
Prompt Engineers & 14\% \\
Translators & 13\% \\
AI Ethics Specialists & 6\% \\
\bottomrule
\end{tabular}
\end{table}


%%%%%%%%%%%%%%%%%%%%%%%%%%%%%%%%%%%%%%%%%%%%%%%%%%%%%%%%%%%%%%%%%%%%%%%%%%%%%%%
\subsection{Risks and Mitigation}
%%%%%%%%%%%%%%%%%%%%%%%%%%%%%%%%%%%%%%%%%%%%%%%%%%%%%%%%%%%%%%%%%%%%%%%%%%%%%%%

\begin{table}[htbp]
\centering
\caption{AI Risks: Negative Consequences Experienced vs. Mitigation Efforts}
\label{tab:vvv1-risks}
\renewcommand{\arraystretch}{1.3}
\begin{tabular}{@{}lcc@{}}
\toprule
\textbf{Risk} & \textbf{Experienced} & \textbf{Mitigating} \\
\midrule
\textbf{Inaccuracy} & \textbf{30\%} & \textbf{54\%} \\
Explainability & 14\% & 28\% \\
Personal/Individual Privacy & 11\% & 38\% \\
Cybersecurity & 10\% & 51\% \\
Regulatory Compliance & 8\% & 43\% \\
IP Infringement & 8\% & 38\% \\
Unauthorized Action & 7\% & 28\% \\
Equity and Fairness & 7\% & 21\% \\
Workforce Displacement & 6\% & 14\% \\
Organizational Reputation & 5\% & 27\% \\
\bottomrule
\end{tabular}
\end{table}

\textbf{Key Finding:} 51\% of organizations using AI have experienced at least one negative consequence. \textbf{High performers report MORE negative consequences}---but also higher mitigation rates. This reflects the learning dynamics of ambitious deployment.


%%%%%%%%%%%%%%%%%%%%%%%%%%%%%%%%%%%%%%%%%%%%%%%%%%%%%%%%%%%%%%%%%%%%%%%%%%%%%%%
\subsection{Integration with EBF: Theoretical Propositions Validated}
\label{sec:vvv1-integration}
%%%%%%%%%%%%%%%%%%%%%%%%%%%%%%%%%%%%%%%%%%%%%%%%%%%%%%%%%%%%%%%%%%%%%%%%%%%%%%%

\begin{tcolorbox}[colback=purple!5!white, colframe=purple!75!black,
    title=Empirical Support for EBF Propositions]

\textbf{Proposition 4.1 (Institutional Complementarity):}
\begin{quote}
Technological adoption requires complementary institutional investments.
\end{quote}
\textit{Evidence:} 62\% still in pilot/experiment phase despite widespread AI availability.

\vspace{0.5em}

\textbf{Proposition 5.3 (Supermodularity):}
\begin{quote}
The value of pursuing one objective increases with pursuit of others.
\end{quote}
\textit{Evidence:} High performers pursue efficiency (84\%) + growth (82\%) + innovation (79\%) simultaneously. Others pursue efficiency (80\%) but only 50\% each for growth/innovation.

\vspace{0.5em}

\textbf{Proposition 6.2 (Context Heterogeneity):}
\begin{quote}
AI complementarities are context-dependent ($\Psi$-modulated).
\end{quote}
\textit{Evidence:} Wide variation in adoption by industry, function, and company size. Technology leads; manufacturing lags. Same technology, different contexts, different outcomes.

\vspace{0.5em}

\textbf{Proposition 7.1 (Leadership as Complementary Asset):}
\begin{quote}
Human capital and organizational commitment are essential complements to technological capability.
\end{quote}
\textit{Evidence:} 3.0× differential in leadership commitment between high performers (48\%) and others (16\%). Leadership is not optional---it's a \textbf{complementary requirement}.

\end{tcolorbox}


%%%%%%%%%%%%%%%%%%%%%%%%%%%%%%%%%%%%%%%%%%%%%%%%%%%%%%%%%%%%%%%%%%%%%%%%%%%%%%%
\subsection{Implications for EBF Productizing (Appendix VVV)}
%%%%%%%%%%%%%%%%%%%%%%%%%%%%%%%%%%%%%%%%%%%%%%%%%%%%%%%%%%%%%%%%%%%%%%%%%%%%%%%

The McKinsey findings directly inform the EBF productizing strategy:

\begin{enumerate}
\item \textbf{Window of Opportunity:} 62\% still in pilot phase = market not yet consolidated
\item \textbf{Agent Adoption Nascent:} Only 23\% scaling agents = EBF Skill timing is right
\item \textbf{High Performers Want Transformation:} 50\% aim for transformative change = demand for comprehensive frameworks like EBF
\item \textbf{Workflow Redesign Critical:} EBF's 10C CORE provides structure for workflow analysis
\item \textbf{Multi-Objective Optimization Needed:} EBF's complementarity analysis addresses the efficiency+growth+innovation requirement
\end{enumerate}


%%%%%%%%%%%%%%%%%%%%%%%%%%%%%%%%%%%%%%%%%%%%%%%%%%%%%%%%%%%%%%%%%%%%%%%%%%%%%%%
\subsection{Methodological Notes}
%%%%%%%%%%%%%%%%%%%%%%%%%%%%%%%%%%%%%%%%%%%%%%%%%%%%%%%%%%%%%%%%%%%%%%%%%%%%%%%

\begin{itemize}
\item \textbf{Survey Period:} June 25 -- July 29, 2025
\item \textbf{Sample Size:} n = 1,993 participants
\item \textbf{Geographic Coverage:} 105 nations
\item \textbf{Weighting:} Adjusted by national contribution to global GDP
\item \textbf{Organization Size:} 38\% of respondents from organizations with $>$\$1B annual revenue
\item \textbf{AI High Performers:} n = 109 (6\% of sample)
\end{itemize}


%%%%%%%%%%%%%%%%%%%%%%%%%%%%%%%%%%%%%%%%%%%%%%%%%%%%%%%%%%%%%%%%%%%%%%%%%%%%%%%
\subsection{Source Reference}
%%%%%%%%%%%%%%%%%%%%%%%%%%%%%%%%%%%%%%%%%%%%%%%%%%%%%%%%%%%%%%%%%%%%%%%%%%%%%%%

\begin{quote}
Singla, A., Sukharevsky, A., Yee, L., Chui, M., Hall, B., \& Balakrishnan, T. (2025). \textit{The state of AI in 2025: Agents, innovation, and transformation}. McKinsey \& Company, QuantumBlack. November 2025.
\end{quote}

The full PDF report is available in the repository as:

\texttt{appendices/the-state-of-ai-2025-agents-innovation\_cmyk-v1.pdf}


%%%%%%%%%%%%%%%%%%%%%%%%%%%%%%%%%%%%%%%%%%%%%%%%%%%%%%%%%%%%%%%%%%%%%%%%%%%%%%%

%%%%%%%%%%%%%%%%%%%%%%%%%%%%%%%%%%%%%%%%%%%%%%%%%%%%%%%%%%%%%%%%%%%%%%%%%%%%%%%
\section{Critical Foundations and Objections}
\label{sec:vvv7-foundations}
%%%%%%%%%%%%%%%%%%%%%%%%%%%%%%%%%%%%%%%%%%%%%%%%%%%%%%%%%%%%%%%%%%%%%%%%%%%%%%%

\subsection{Objection 1: Scope Limitations}

\textbf{Objection:} The framework presented may have limited applicability
outside the specific domain contexts discussed.

\textbf{Response:} We acknowledge domain-specificity. The framework provides
general principles that should be calibrated to specific contexts. Cross-domain
validation remains an open research question.

\subsection{Objection 2: Measurement Challenges}

\textbf{Objection:} Key constructs may be difficult to measure reliably in
practice.

\textbf{Response:} We recommend using validated scales and instruments where
available, and developing domain-specific proxies where necessary. Sensitivity
analysis should assess robustness to measurement uncertainty.



%%%%%%%%%%%%%%%%%%%%%%%%%%%%%%%%%%%%%%%%%%%%%%%%%%%%%%%%%%%%%%%%%%%%%%%%%%%%%%%


%%%%%%%%%%%%%%%%%%%%%%%%%%%%%%%%%%%%%%%%%%%%%%%%%%%%%%%%%%%%%%%%%%%%%%%%%%%%%%%
\section{Summary}
\label{sec:vvv7-summary}
%%%%%%%%%%%%%%%%%%%%%%%%%%%%%%%%%%%%%%%%%%%%%%%%%%%%%%%%%%%%%%%%%%%%%%%%%%%%%%%

\begin{tcolorbox}[colback=gray!5!white, colframe=gray!75!black,
    title=Appendix VVV7 Summary]

\textbf{Core Insight:}
This appendix provides key analysis and findings relevant to the EBF framework.

\textbf{Key Contributions:}
\begin{itemize}[nosep]
    \item Theoretical foundations and empirical evidence
    \item Parameter estimates and model validation
    \item Integration with 10C CORE architecture
\end{itemize}

\textbf{Framework Integration:}
The findings connect to WHO, WHAT, HOW, WHEN, WHERE, AWARE, READY, and STAGE dimensions.

\end{tcolorbox}

\section{Glossary of Symbols}
\label{sec:vvv7-glossary}
%%%%%%%%%%%%%%%%%%%%%%%%%%%%%%%%%%%%%%%%%%%%%%%%%%%%%%%%%%%%%%%%%%%%%%%%%%%%%%%

For comprehensive definitions, see Appendix G (Master Glossary).

\begin{table}[htbp]
\centering
\caption{Key Symbols in Appendix VVV7}
\begin{tabular}{lll}
\toprule
\textbf{Symbol} & \textbf{Meaning} & \textbf{Reference} \\
\midrule
$\gamma$ & Complementarity parameter & Appendix B \\
$\Psi$ & Context vector & Appendix V \\
$U$ & Utility function & Chapter 3 \\
\bottomrule
\end{tabular}
\end{table}



%%%%%%%%%%%%%%%%%%%%%%%%%%%%%%%%%%%%%%%%%%%%%%%%%%%%%%%%%%%%%%%%%%%%%%%%%%%%%%%
\section{Open Issues and Research Agenda}
\label{sec:x-open}
%%%%%%%%%%%%%%%%%%%%%%%%%%%%%%%%%%%%%%%%%%%%%%%%%%%%%%%%%%%%%%%%%%%%%%%%%%%%%%%

\begin{enumerate}
    \item \textbf{Empirical Validation:} Further testing across diverse contexts
    \item \textbf{Parameter Refinement:} Improved estimation methods needed
    \item \textbf{Integration:} Enhanced cross-appendix linkages
\end{enumerate}

\section{References}
\label{sec:references}
%%%%%%%%%%%%%%%%%%%%%%%%%%%%%%%%%%%%%%%%%%%%%%%%%%%%%%%%%%%%%%%%%%%%%%%%%%%%%%%

See master bibliography (\texttt{bcm\_master.bib}) for complete citations.

\nocite{bcm_master}

% =============================================================================
% END OF VVV1
% =============================================================================
